% ===== PRÉAMBULE CANONIQUE — à coller VERBATIM en tête de CHAQUE .tex =====
\documentclass[11pt,a4paper]{article}
\usepackage[utf8]{inputenc}\usepackage[T1]{fontenc}\usepackage{lmodern}
\usepackage[french]{babel}
\usepackage{amsmath,amssymb,mathtools}\usepackage{icomma}
\usepackage{siunitx}\sisetup{locale=FR}  % \qty{}{}, \si{}, \ang{}
\usepackage{xcolor}\usepackage{colortbl}
\usepackage{tikz}\usepackage{pgfplots}\pgfplotsset{compat=1.18}
\usepackage[siunitx,europeanresistors]{circuitikz} % circuits (élec.)
\usepackage[most]{tcolorbox}\usepackage{enumitem}\usepackage{array,booktabs}
\usepackage[a4paper,margin=2.2cm]{geometry}
\usetikzlibrary{arrows.meta,calc,angles,quotes,positioning,patterns,%
  backgrounds,shapes.geometric,decorations.pathmorphing,decorations.markings,intersections}
\definecolor{primary}{RGB}{222,49,99}\definecolor{primarydark}{RGB}{150,22,55}
\definecolor{defcol}{RGB}{0,114,178}\definecolor{methcol}{RGB}{52,92,148}
\definecolor{excol}{RGB}{0,135,90}\definecolor{attncol}{RGB}{200,40,55}
\tcbset{boxbase/.style={enhanced,breakable,boxrule=0.9pt,arc=2.5pt,
  left=3mm,right=3mm,top=2.3mm,bottom=2.3mm,fonttitle=\bfseries,coltitle=white}}
% --- boîtes TUTORIEL (noms EXACTS, ne pas renommer) ---
\newtcolorbox{cle}[1][]{boxbase,colback=defcol!6,colframe=defcol,
  title={\textsc{À retenir}\ifx\relax#1\relax\relax\else\textmd{ : \itshape #1}\fi}}
\newtcolorbox{methode}[1][]{boxbase,colback=methcol!7,colframe=methcol,sharp corners,
  title={\textsc{Méthode}\ifx\relax#1\relax\relax\else\textmd{ --- \itshape #1}\fi}}
\newtcolorbox{exemplebox}[1][]{boxbase,colback=excol!5,colframe=excol,
  title={\textsc{Exemple}\ifx\relax#1\relax\relax\else\textmd{ : \itshape #1}\fi}}
\newtcolorbox{piege}[1][]{boxbase,colback=attncol!6,colframe=attncol,
  title={\textsc{Piège}\ifx\relax#1\relax\relax\else\textmd{ : \itshape #1}\fi}}
\newtcolorbox{remarque}[1][]{boxbase,colback=black!4,colframe=black!45,
  title={\textsc{Remarque}\ifx\relax#1\relax\relax\else\textmd{ : \itshape #1}\fi}}
% --- boîtes EXERCICE / CORRIGÉ (noms EXACTS, auto-comptées) ---
\newtcolorbox[auto counter]{exo}[1][]{boxbase,colback=primary!4,colframe=primary,
  title={\textsc{Exercice}~\thetcbcounter\ifx\relax#1\relax\relax\else\textmd{ --- \itshape #1}\fi}}
\newtcolorbox[auto counter]{corrige}[1][]{boxbase,colback=excol!4,colframe=excol,
  title={\textsc{Corrigé}~\thetcbcounter\ifx\relax#1\relax\relax\else\textmd{ --- \itshape #1}\fi}}
\newcommand{\vv}[1]{\vec{#1}}\newcommand{\ds}{\displaystyle}
\newcommand{\R}{\mathbb{R}}\newcommand{\N}{\mathbb{N}}\newcommand{\Z}{\mathbb{Z}}
% (un \usepackage{fancyhdr}+en-tête est possible ; il est ignoré côté web)
% =========================================================================
\begin{document}
\begin{corrige}[sous une ligne à haute tension]
\begin{enumerate}[label=\alph*)]
  \item $B=\dfrac{2\cdot10^{-7}\times100}{0{,}4}=\dfrac{2\cdot10^{-5}}{0{,}4}=\qty{5e-5}{\tesla}$
        (autant que le champ terrestre).
  \item $B=\dfrac{2\cdot10^{-7}\times100}{30}=\dfrac{2\cdot10^{-5}}{30}\approx\qty{6.7e-7}{\tesla}$.
        Rapport : $\dfrac{B_T}{B}=\dfrac{5\cdot10^{-5}}{6{,}7\cdot10^{-7}}\approx 75$ : le champ de la
        ligne est \textbf{environ 75 fois plus faible} que le champ terrestre. C'est pourquoi on place
        les câbles très haut : la décroissance en $1/d$ suffit à rendre le champ négligeable au sol.
\end{enumerate}
\end{corrige}

\begin{corrige}[mettre le solénoïde à l'échelle]
\begin{enumerate}[label=\alph*)]
  \item $B=\dfrac{\mu_0 N I}{L}$ : $B$ est proportionnel à $N$, à $I$, et à $1/L$. Donc
        \[ \text{facteur}= \underbrace{2}_{N\times2}\times\underbrace{3}_{I\times3}\times
           \underbrace{2}_{L\div2\ \Rightarrow\ 1/L\times2}=\mathbf{12}. \]
  \item $B=12\times B_0=12\times2\cdot10^{-3}=\qty{2.4e-2}{\tesla}=\qty{24}{\milli\tesla}.$
\end{enumerate}
\end{corrige}

\begin{corrige}[combien de spires pour la bobine ?]
On isole $N$ dans $B=\dfrac{N\,\mu_0 I}{2r}$, avec $r=\qty{0.03}{\metre}$ :
\[ N=\frac{B\,(2r)}{\mu_0 I}=\frac{2\cdot10^{-3}\times0{,}06}{1{,}26\cdot10^{-6}\times4}
    =\frac{1{,}2\cdot10^{-4}}{5{,}03\cdot10^{-6}}\approx 24\ \text{spires}. \]
\end{corrige}

\begin{corrige}[deux fils : Biot-Savart puis superposition]
$d=\qty{0.03}{\metre}$ pour chaque fil.
\begin{enumerate}[label=\alph*)]
  \item $B_1=\dfrac{2\cdot10^{-7}\times10}{0{,}03}=\dfrac{2\cdot10^{-6}}{0{,}03}\approx\qty{6.7e-5}{\tesla}$ ;
        \quad $B_2=\dfrac{2\cdot10^{-7}\times20}{0{,}03}=\dfrac{4\cdot10^{-6}}{0{,}03}\approx\qty{1.3e-4}{\tesla}$.
  \item Les deux courants sortent : main droite $\Rightarrow$ en $M$, $\vv{B}_1$ (dû au fil de gauche)
        est dirigé \textbf{vers le haut}, $\vv{B}_2$ (dû au fil de droite) \textbf{vers le bas}. Ils
        sont \textbf{opposés}, donc on soustrait :
        \[ B_{\text{tot}}=B_2-B_1=1{,}3\cdot10^{-4}-6{,}7\cdot10^{-5}\approx\qty{6.7e-5}{\tesla}, \]
        dirigé \textbf{vers le bas} (sens de $\vv{B}_2$, le plus intense).
\end{enumerate}
\end{corrige}

\end{document}
